\documentclass[a4paper,12pt]{article}

% Paquetes esenciales
\usepackage[utf8]{inputenc}
\usepackage[T1]{fontenc}
\usepackage[spanish]{babel}
\usepackage{geometry}
\usepackage{graphicx}
\usepackage{hyperref}
\usepackage{xcolor}
\usepackage{tikz}
\usetikzlibrary{shapes, arrows, positioning, calc, shadows.blur}
\usepackage{float}
\usepackage{titlesec}
\usepackage{fancyhdr}
\usepackage{listings}
\usepackage{courier}

% Configuración de página
\geometry{margin=1in}

% Colores corporativos (basados en el tema DeFi)
\definecolor{defiblue}{RGB}{30, 58, 138}
\definecolor{defilightblue}{RGB}{59, 130, 246}
\definecolor{defigreen}{RGB}{16, 185, 129}
\definecolor{defidark}{RGB}{17, 24, 39}

% Configuración de hipervínculos
\hypersetup{
    colorlinks=true,
    linkcolor=defiblue,
    filecolor=magenta,      
    urlcolor=defiblue,
    citecolor=defigreen,
}

% Configuración de código
\lstset{
    language=HTML,
    basicstyle=\ttfamily\small,
    backgroundcolor=\color{gray!10},
    frame=single,
    breaklines=true,
    postbreak=\mbox{\textcolor{red}{$\hookrightarrow$}\space},
}

% Encabezado y pie de página
\pagestyle{fancy}
\fancyhf{}
\rhead{\textbf{DefiLab}}
\lhead{Documentación Técnica}
\rfoot{Página \thepage}

% Título
\title{\textbf{\Huge DefiLab}\\[0.5cm]
\Large Biblioteca de Componentes UI y Dashboard Admin\\[0.3cm]
\large Documentación Técnica y Arquitectura}
\author{Equipo de Desarrollo DefiLab}
\date{\today}

\begin{document}

\maketitle

\begin{abstract}
\noindent Este documento proporciona una descripción técnica completa de \textbf{DefiLab}, una biblioteca de componentes UI y plantilla de administración moderna diseñada para aplicaciones DeFi y dashboards empresariales. Se incluye la arquitectura del proyecto, diagramas de flujo, estructura de componentes y guías de implementación.
\end{abstract}

\tableofcontents
\newpage

% ============================================================================
% SECCIÓN 1: INTRODUCCIÓN
% ============================================================================
\section{Introducción}

\subsection{Propósito del Proyecto}
DefiLab es una solución integral para el desarrollo rápido de interfaces web financieras. Proporciona componentes preconstruidos, estilos responsivos y plantillas de dashboard que aceleran el tiempo de desarrollo en un 60\%.

\subsection{Características Principales}
\begin{itemize}
    \item \textbf{Visualización de Datos}: Gráficos interactivos (líneas, barras, tortas)
    \item \textbf{Componentes Rich-UI}: Botones, tarjetas, modales, alertas
    \item \textbf{Formularios Avanzados}: Validaciones en tiempo real
    \item \textbf{Autenticación}: Plantillas completas de login/registro
    \item \textbf{Diseño Responsivo}: Mobile-first compatible con todos los dispositivos
    \item \textbf{Iconografía}: Soporte para FontAwesome y Themify Icons
\end{itemize}

% ============================================================================
% SECCIÓN 2: ARQUITECTURA DEL SISTEMA
% ============================================================================
\section{Arquitectura del Sistema}

\subsection{Diagrama de Arquitectura General}
\begin{figure}[H]
    \centering
    \begin{tikzpicture}[
        node distance=1.5cm,
        box/.style={rectangle, draw, fill=white, blur shadow, minimum width=3cm, minimum height=1cm, rounded corners},
        arrow/.style={->, >=stealth, thick, defiblue}
    ]
        % Nodos
        \node[box, fill=defiblue!10] (user) {Usuario Final};
        \node[box, below=of user, fill=defilightblue!10] (browser) {Navegador Web};
        \node[box, below=of browser, fill=defigreen!10] (html) {index.html};
        \node[box, left=of html, fill=defigreen!10] (css) {styles/css/};
        \node[box, right=of html, fill=defigreen!10] (js) {styles/js/};
        \node[box, below=of html, fill=defiblue!10] (components) {Componentes UI};
        
        % Flechas
        \draw[arrow] (user) -- (browser);
        \draw[arrow] (browser) -- (html);
        \draw[arrow] (html) -- (css);
        \draw[arrow] (html) -- (js);
        \draw[arrow] (html) -- (components);
    \end{tikzpicture}
    \caption{Arquitectura general de DefiLab}
    \label{fig:arquitectura}
\end{figure}

\subsection{Estructura de Directorios}
\begin{lstlisting}[language=bash, caption=Estructura del proyecto]
/workspace
├── styles/               # Directorio principal
│   ├── css/              # Hojas de estilo CSS
│   │   ├── style.css     # Estilos principales
│   │   └── responsive.css # Media queries
│   ├── js/               # Scripts JavaScript
│   │   ├── main.js       # Lógica principal
│   │   └── charts.js     # Configuración de gráficos
│   ├── img/              # Recursos de imagen
│   │   └── logo.png      # Logotipo del proyecto
│   └── index.html        # Punto de entrada principal
├── docs/                 # Documentación (este archivo)
└── README.md             # Documentación Markdown
\end{lstlisting}

% ============================================================================
% SECCIÓN 3: DIAGRAMAS DE FLUJO
% ============================================================================
\section{Diagramas de Flujo}

\subsection{Flujo de Carga de la Aplicación}
\begin{figure}[H]
    \centering
    \begin{tikzpicture}[
        node distance=1.2cm,
        startstop/.style={rectangle, rounded corners, minimum width=2cm, minimum height=0.8cm, text centered, draw=black, fill=red!30},
        process/.style={rectangle, minimum width=2cm, minimum height=0.8cm, text centered, draw=black, fill=defiblue!20},
        decision/.style={diamond, minimum width=2cm, minimum height=0.8cm, text centered, draw=black, fill=defilightblue!20},
        arrow/.style={->, >=stealth, thick}
    ]
        % Nodos
        \node[startstop] (start) {Inicio};
        \node[process, below=of start] (load-html) {Cargar index.html};
        \node[process, below=of load-html] (load-css) {Cargar CSS};
        \node[process, below=of load-css] (load-js) {Cargar JS};
        \node[decision, below=of load-js] (check-dom) {¿DOM listo?};
        \node[process, below=of check-dom] (init-components) {Inicializar componentes};
        \node[process, below=of init-components] (render-charts) {Renderizar gráficos};
        \node[startstop, below=of render-charts] (end) {Aplicación lista};
        
        % Flechas
        \draw[arrow] (start) -- (load-html);
        \draw[arrow] (load-html) -- (load-css);
        \draw[arrow] (load-css) -- (load-js);
        \draw[arrow] (load-js) -- (check-dom);
        \draw[arrow] (check-dom) -- node[right] {Sí} (init-components);
        \draw[arrow] (check-dom.west) -- ++(-1,0) |- node[left] {No} (load-js);
        \draw[arrow] (init-components) -- (render-charts);
        \draw[arrow] (render-charts) -- (end);
    \end{tikzpicture}
    \caption{Flujo de carga de la aplicación}
    \label{fig:flujo-carga}
\end{figure}

\subsection{Flujo de Autenticación de Usuario}
\begin{figure}[H]
    \centering
    \begin{tikzpicture}[
        node distance=1.2cm,
        startstop/.style={rectangle, rounded corners, minimum width=2cm, minimum height=0.8cm, text centered, draw=black, fill=red!30},
        process/.style={rectangle, minimum width=2cm, minimum height=0.8cm, text centered, draw=black, fill=defiblue!20},
        decision/.style={diamond, minimum width=2cm, minimum height=0.8cm, text centered, draw=black, fill=defilightblue!20},
        arrow/.style={->, >=stealth, thick}
    ]
        % Nodos
        \node[startstop] (start) {Usuario ingresa};
        \node[process, below=of start] (login-form) {Completa formulario login};
        \node[process, below=of login-form] (submit) {Enviar credenciales};
        \node[decision, below=of submit] (validate) {¿Credenciales válidas?};
        \node[process, right=of validate, fill=defigreen!20] (dashboard) {Redirigir al Dashboard};
        \node[process, left=of validate, fill=red!20] (error) {Mostrar error};
        \node[startstop, below=of validate] (end-success) {Sesión iniciada};
        
        % Flechas
        \draw[arrow] (start) -- (login-form);
        \draw[arrow] (login-form) -- (submit);
        \draw[arrow] (submit) -- (validate);
        \draw[arrow] (validate) -- node[above] {Sí} (dashboard);
        \draw[arrow] (validate) -- node[above] {No} (error);
        \draw[arrow] (error.west) -- ++(-0.5,0) |- (login-form);
        \draw[arrow] (dashboard) -- (end-success);
    \end{tikzpicture}
    \caption{Flujo de autenticación de usuario}
    \label{fig:flujo-auth}
\end{figure}

\subsection{Flujo de Renderizado de Gráficos}
\begin{figure}[H]
    \centering
    \begin{tikzpicture}[
        node distance=1.2cm,
        startstop/.style={rectangle, rounded corners, minimum width=2cm, minimum height=0.8cm, text centered, draw=black, fill=red!30},
        process/.style={rectangle, minimum width=2cm, minimum height=0.8cm, text centered, draw=black, fill=defiblue!20},
        decision/.style={diamond, minimum width=2cm, minimum height=0.8cm, text centered, draw=black, fill=defilightblue!20},
        arrow/.style={->, >=stealth, thick}
    ]
        % Nodos
        \node[startstop] (start) {Inicio renderizado};
        \node[process, below=of start] (fetch-data) {Obtener datos};
        \node[decision, below=of fetch-data] (has-data) {¿Hay datos?};
        \node[process, below=of has-data] (config-chart) {Configurar gráfico};
        \node[process, below=of config-chart] (draw-chart) {Dibujar gráfico};
        \node[process, below=of draw-chart] (animate) {Aplicar animaciones};
        \node[startstop, below=of animate] (end) {Gráfico listo};
        \node[process, left=of has-data, fill=yellow!30] (fallback) {Usar datos mock};
        
        % Flechas
        \draw[arrow] (start) -- (fetch-data);
        \draw[arrow] (fetch-data) -- (has-data);
        \draw[arrow] (has-data) -- node[right] {Sí} (config-chart);
        \draw[arrow] (has-data.west) -- node[above] {No} (fallback);
        \draw[arrow] (fallback.north) -- ++(0,0.5) -| (fetch-data);
        \draw[arrow] (config-chart) -- (draw-chart);
        \draw[arrow] (draw-chart) -- (animate);
        \draw[arrow] (animate) -- (end);
    \end{tikzpicture}
    \caption{Flujo de renderizado de gráficos}
    \label{fig:flujo-charts}
\end{figure}

% ============================================================================
% SECCIÓN 4: COMPONENTES UI
% ============================================================================
\section{Componentes UI Disponibles}

\subsection{Tabla de Componentes}
\begin{table}[H]
    \centering
    \begin{tabular}{|l|l|l|}
        \hline
        \textbf{Categoría} & \textbf{Componentes} & \textbf{Ubicación} \\ \hline
        Gráficos & Barras, Líneas, Tortas & styles/js/charts.js \\ \hline
        Formularios & Inputs, Selects, Validaciones & styles/css/forms.css \\ \hline
        Tablas & Básicas, Data Tables & styles/css/tables.css \\ \hline
        Navegación & Menús, Tabs, Breadcrumbs & styles/css/nav.css \\ \hline
        Autenticación & Login, Registro, Recovery & styles/index.html \\ \hline
        Utilidades & Alertas, Modales, Tooltips & styles/css/utils.css \\ \hline
    \end{tabular}
    \caption{Resumen de componentes disponibles}
    \label{tab:componentes}
\end{table}

\subsection{Ejemplo de Uso: Gráfico de Líneas}
\begin{lstlisting}[caption=Ejemplo de implementación de gráfico]
<!-- HTML -->
<div class="chart-container">
    <canvas id="lineChart"></canvas>
</div>

<!-- JavaScript -->
<script>
    const ctx = document.getElementById('lineChart').getContext('2d');
    const chart = new Chart(ctx, {
        type: 'line',
        data: {
            labels: ['Ene', 'Feb', 'Mar', 'Abr'],
            datasets: [{
                label: 'Precio USD',
                data: [100, 120, 110, 130],
                borderColor: '#3b82f6'
            }]
        }
    });
</script>
\end{lstlisting}

% ============================================================================
% SECCIÓN 5: GUÍA DE IMPLEMENTACIÓN
% ============================================================================
\section{Guía de Implementación}

\subsection{Requisitos Previos}
\begin{itemize}
    \item Navegador web moderno (Chrome, Firefox, Safari, Edge)
    \item Servidor web local (opcional, pero recomendado)
    \item Conocimientos básicos de HTML/CSS/JS
\end{itemize}

\subsection{Pasos de Instalación}
\begin{enumerate}
    \item Clonar el repositorio:
    \begin{lstlisting}[language=bash]
git clone https://github.com/tu-usuario/defilab.git
cd defilab
    \end{lstlisting}
    
    \item Abrir el archivo principal:
    \begin{lstlisting}[language=bash]
# Opción A: Abrir directamente en el navegador
open styles/index.html

# Opción B: Usar un servidor local
python -m http.server 8000
# Luego visitar http://localhost:8000/styles/index.html
    \end{lstlisting}
    
    \item Personalizar componentes según necesidades del proyecto.
\end{enumerate}

% ============================================================================
% SECCIÓN 6: DIAGRAMA DE CLASES (COMPONENTES)
% ============================================================================
\section{Diagrama de Componentes}

\begin{figure}[H]
    \centering
    \begin{tikzpicture}[
        node distance=1.5cm,
        component/.style={rectangle, draw, fill=white, blur shadow, minimum width=2.5cm, minimum height=1cm, rounded corners},
        arrow/.style={->, >=stealth, thick, defiblue}
    ]
        % Nodos
        \node[component, fill=defiblue!20] (core) {Core System};
        \node[component, above left=of core] (layout) {Layout Manager};
        \node[component, above right=of core] (theme) {Theme Engine};
        \node[component, below left=of core] (ui) {UI Components};
        \node[component, below right=of core] (data) {Data Visualization};
        \node[component, left=of ui] (forms) {Forms};
        \node[component, right=of ui] (auth) {Auth};
        \node[component, left=of data] (charts) {Charts};
        \node[component, right=of data] (maps) {Maps};
        
        % Flechas
        \draw[arrow] (layout) -- (core);
        \draw[arrow] (theme) -- (core);
        \draw[arrow] (core) -- (ui);
        \draw[arrow] (core) -- (data);
        \draw[arrow] (ui) -- (forms);
        \draw[arrow] (ui) -- (auth);
        \draw[arrow] (data) -- (charts);
        \draw[arrow] (data) -- (maps);
    \end{tikzpicture}
    \caption{Diagrama de componentes del sistema}
    \label{fig:componentes}
\end{figure}

% ============================================================================
% SECCIÓN 7: CONTRIBUCIÓN
% ============================================================================
\section{Contribución}

Para contribuir al proyecto, seguir el siguiente flujo:

\begin{figure}[H]
    \centering
    \begin{tikzpicture}[
        node distance=1cm,
        process/.style={rectangle, minimum width=2cm, minimum height=0.8cm, text centered, draw=black, fill=defiblue!20},
        arrow/.style={->, >=stealth, thick}
    ]
        % Nodos
        \node[process] (fork) {Fork del repo};
        \node[process, below=of fork] (clone) {Clonar fork};
        \node[process, below=of clone] (branch) {Crear rama feature};
        \node[process, below=of branch] (commit) {Commit cambios};
        \node[process, below=of commit] (push) {Push a GitHub};
        \node[process, below=of push] (pr) {Crear Pull Request};
        
        % Flechas
        \draw[arrow] (fork) -- (clone);
        \draw[arrow] (clone) -- (branch);
        \draw[arrow] (branch) -- (commit);
        \draw[arrow] (commit) -- (push);
        \draw[arrow] (push) -- (pr);
    \end{tikzpicture}
    \caption{Flujo de contribución al proyecto}
    \label{fig:contribucion}
\end{figure}

% ============================================================================
% SECCIÓN 8: CONCLUSIONES
% ============================================================================
\section{Conclusiones}

DefiLab proporciona una base sólida para el desarrollo de aplicaciones DeFi y dashboards empresariales. Su arquitectura modular permite una fácil personalización y extensión. Los diagramas presentados en este documento facilitan la comprensión del flujo de trabajo y la estructura del sistema.

\subsection*{Próximos Pasos}
\begin{itemize}
    \item Implementar soporte para temas oscuros
    \item Añadir más tipos de gráficos
    \item Integrar con APIs de criptomonedas en tiempo real
    \item Mejorar la accesibilidad (WCAG)
\end{itemize}

% ============================================================================
% BIBLIOGRAFÍA
% ============================================================================
\begin{thebibliography}{9}
    \bibitem{html5} MDN Web Docs. \textit{HTML5 Reference}. Mozilla Developer Network.
    \bibitem{css3} MDN Web Docs. \textit{CSS3 Reference}. Mozilla Developer Network.
    \bibitem{chartjs} Chart.js Documentation. \textit{Getting Started}. chartjs.org
    \bibitem{bootstrap} Bootstrap Documentation. \textit{Components}. getbootstrap.com
\end{thebibliography}

\end{document}
